\begin{textsample}{NEG Dim 2 – 2015\_c5 – Score -1.00 – 2015\_c5/enade\_2015\_ci\_cont\_383e7cacfb.txt}  \label{ex:f2_neg_155}
CAPÍTULO 5 DISTRIBUIÇÃO DOS CONCEITOS 5.1 PANORAMA NACIONAL DA DISTRIBUIÇÃO DOS CONCEITOS O Conceito Enade é um indicador de qualidade que avalia o desempenho dos estudantes a partir dos resultados do Enade.
Os valores possíveis vão de 1 ( pior situação ) a 5 ( melhor situação ).
As Unidades de Observação com apenas um ou sem nenhum concluinte participante não obtêm o Conceito Enade, ficando Sem Conceito ( SC ).
A Tabela 5.1 apresenta a quantidade e distribuição de cursos de Ciências Contábeis participantes do Enade/2015, por faixa de conceito e Grande Região.
Mais uma vez ressaltando que a diferença entre os cursos tabulados neste capítulo e no capítulo 2 corresponde aos cursos sem conceito, em princípio, aqueles sem alunos concluintes que participassem da prova ou que tivessem tão somente um aluno concluinte18.
Observando -se os dados da Tabela 5.1, nota -se que, dos 1.044 cursos participantes, 431 ( 41,3 \% ) classificaram -se com conceito 3, o valor modal.
Este foi também o conceito modal em mais três regiões : Nordeste ( 41,6 \% ), Sudeste ( 40,5 \% ) e Sul ( 48,2 \% ).
Nas regiões Norte ( 40,8 \% ) e Centro-Oeste ( 40,9 \% ), a moda foi o conceito 2.
O conceito 2 foi o segundo mais frequente em nível nacional ( 32,1 \%, correspondendo a 335 cursos ) e o conceito 4, o terceiro ( 18,0 \%, correspondendo a 188 cursos ).
Houve, ainda, 43 cursos ( 4,1 \% ) que receberam conceito 5 e 31 cursos ( 3,0 \% ) que receberam conceito 1.
Dos 1.044 cursos de Ciências Contábeis, 16 cursos ( 1,5 \% ) ficaram sem conceito ( SC ). 18 Estes cursos com somente um concluinte não foram considerados no capítulo 2.
A região Norte participou com 71 cursos ou 6,8 \% do total nacional.
Nesta região, o conceito 2 foi atribuído a 29 cursos dos 71 participantes, o que equivale a 40,8 \% do total regional, além de ser o conceito modal na região.
Apenas um curso ( 1,4 \% ) ficou sem conceito ( SC ), nenhum curso teve conceito 5 e cinco cursos ( 7,0 \% ) foram avaliados com conceito 1.
Dos cursos restantes, 24 cursos ( 33,8 \% ) obtiveram conceito 3 e 12 cursos ( 16,9 \% ) tiveram conceito 4.
A região Norte participou com 71 cursos ou 6,8 \% do total nacional.
Nesta região, o conceito 2 foi atribuído a 29 cursos dos 71 participantes, o que equivale a 40,8 \% do total regional, além de ser o conceito modal na região.
Apenas um curso ( 1,4 \% ) ficou sem conceito ( SC ), nenhum curso teve conceito 5 e cinco cursos ( 7,0 \% ) foram avaliados com conceito 1.
Dos cursos restantes, 24 cursos ( 33,8 \% ) obtiveram conceito 3 e 12 cursos ( 16,9 \% ) tiveram conceito 4.
A região Nordeste participou com 190 cursos ou 18,2 \% do total nacional.
Nesta região, 79 cursos, 41,6 \% em termos regionais, obtiveram conceito 3, o conceito modal para a região.
Os conceitos 1 e 5 foram atribuídos a, respectivamente, nove e três cursos ( 4,7 \% e 1,6 \% ).
Dos cursos restantes, 70 ( 36,8 \% ) obtiveram o conceito 2, 27 cursos ( 14,2 \% ) obtiveram o conceito 4 e apenas dois cursos ficaram sem conceito ( SC ), correspondendo a 1,1 \% nessa região.
Dos 444 cursos participantes da região Sudeste ( 42,5 \% do total nacional ), o conceito 3 também foi o conceito modal, tendo sido obtido por 180 cursos ( 40,5 \% ).
O conceito 1 foi atribuído a nove cursos ( 2,0 \% ), o conceito 5 foi atribuído a 29 cursos ( 6,5 \% ) e oito cursos ( 1,8 \% ) ficaram sem conceito ( SC ).
O conceito 2 foi obtido por 133 cursos ( 30,0 \% ), e obtiveram conceito 4 os 85 cursos restantes ( 19,1 \% ).
Os 224 cursos da região Sul ficaram \textbf{distribuídos} em todas as faixas de conceitos, exceto na faixa de conceito 1, que não foi atribuído a nenhum curso.
Há predominância, também nesta região, do conceito 3 com 48,2 \% ( conceito modal ), ou seja, 108 dos 224 cursos participantes na região Sul.
O conceito 2 foi atribuído a 56 cursos ( 25,0 \% ) e o conceito 4, a 51 cursos ( 22,8 \% ).
O conceito 5 foi recebido por seis cursos ( 2,7 \% ).
Apenas três cursos na região Sul ficaram sem conceito ( SC ).
Já dos 115 cursos participantes na região Centro-Oeste, 47 receberam conceito 2, o conceito modal ( 40,9 \% ).
Nesta região, dois cursos ficaram sem conceito ( SC ) e oito cursos ( 7,0 \% ) obtiveram conceito 1.
Já 40 cursos foram alocados no conceito 3, correspondendo a 34,8 \% dos cursos nessa região.
Cinco cursos foram avaliados com conceito 5 ( 4,3 \% ), e 13 cursos ficaram com conceito 4, correspondendo a 11,3 \%, nessa região.
O Gráfico 5.1 apresenta a distribuição do Conceito Enade segundo Grande Região : poligonais mais à esquerda representam Grandes Regiões com pior distribuição, e poligonais mais à direita, Grandes Regiões com melhores conceitos.
As regiões Norte ( linha azul ) e Centro-Oeste ( linha preta ) apresentam os cursos em IES com conceitos mais baixos, ao passo que as regiões Sudeste ( linha vermelha ) e Sul ( linha roxa ) apresentam as distribuições com valores maiores.
Cursos em IES da região Nordeste ( linha verde ) apresenta uma situação intermediária. 5.2 CONCEITOS POR CATEGORIA ADMINISTRATIVA E POR GRANDE REGIÃO A Tabela 5.2 apresenta a distribuição dos cursos participantes do Enade/2015 de Ciências Contábeis, por Categoria Administrativa, de acordo com os conceitos por eles alcançados, segundo as Grandes Regiões brasileiras.
Dos 1.044 cursos participantes, 148 ( 14,2 \% ) eram ministrados em Instituições Públicas e 896 ( 85,8 \% ), em Instituições Privadas.
De acordo com as informações da Tabela 5.2, em termos nacionais, entre Instituições Públicas e Privadas nota -se que das 43 IES que receberam o conceito 5, 16 eram Públicas e 27, Privadas.
Dos 148 cursos participantes de IES Públicas, o conceito 4 foi o valor modal, atribuído a 58 cursos ( 39,2 \% ).
Entre os demais cursos de Instituições Públicas participantes, seis cursos ( 4,1 \% ) foram avaliados com conceito 1, 25 cursos obtiveram conceito 2 ( 16,9 \% da categoria ), 41 cursos receberam conceito 3 ( 27,7 \% ) e, como já comentado, 16 foram alocados ao conceito 5 ( 10,8 \% ).
Nesta categoria, somente dois cursos ( 1,4 \% ) ficaram sem conceito ( SC ).
Na rede privada, o conceito modal foi 3, com 390 cursos dos 896 da categoria.
Entre os demais cursos participantes, 25 receberam conceito 1, e 310, conceito 2.
O conceito 4 foi atribuído a 130 cursos e o conceito 5, a 27 cursos.
Nesta Categoria Administrativa, outros 14 cursos ficaram sem conceito ( SC ).
Na análise por região, observa -se que, na região Norte, as Instituições Privadas participaram com 58 cursos ( 81,7 \% do total regional ), dos quais apenas um curso ( 1,7 \% ) ficou sem conceito ( SC ), dois cursos ( 3,4 \% ) com conceito 1 e nenhum curso obteve conceito 5.
O conceito modal para as Instituições Privadas na região foi 2, com 26 cursos, correspondendo a 44,8 \%.
Quanto aos demais conceitos, 22 cursos obtiveram conceito 3 e sete cursos tiraram conceito 4.
As Instituições Públicas participaram com 13 cursos na região Norte ( 18,3 \% do total regional ), dos quais nenhum ficou sem conceito ( SC ) nem obteve conceito 5.
Apenas dois cursos ( 15,4 \% ) obtiveram conceito 3 e cinco cursos ( 38,5 \% ) obtiveram conceito 4, o valor modal.
No que diz respeito aos demais conceitos, três cursos ( 23,1 \% ) obtiveram conceito 1 e três cursos obtiveram conceito 2.
Na região Nordeste, a rede privada concentrou 146 dos 190 cursos participantes, o equivalente a 76,8 \% do total da região.
Dentre os cursos oferecidos por IES Privadas no Nordeste, dois cursos ficaram sem conceito ( SC ) e outros dois cursos foram alocados no conceito 5.
Outros 61 cursos obtiveram conceito 2 ( 41,8 \% ), o valor modal e 59 cursos receberam conceito 3 ( 40,4 \% ).
Sobre os demais cursos, oito cursos ( 5,5 \% ) tiraram conceito 1, e 14 cursos ( 9,6 \% ) obtiveram conceito 4.
As Instituições Públicas dessa região participaram com 44 cursos ( 23,2 \% ), dos quais 20 ( 45,5 \% ) obtiveram conceito 3, o valor modal, seguido por 13 cursos que receberam conceito 4 ( 29,5 \% ).
Um curso ( 2,3 \% ) foi alocado no conceito 1, enquanto outros nove obtiveram conceito 2.
Nenhum curso ficou sem conceito ( SC ) nesta combinação de categoria e região e apenas 1 curso ( 2,3 \% ) obteve o conceito 5.
Na região Sudeste, a proporção de cursos da rede privada, 92,6 \%, foi mais elevada do que nas demais regiões brasileiras, correspondendo a 411 dos 444 cursos participantes.
Nesta categoria, na região Sudeste, o conceito modal foi 3 ( 176 cursos ).
Nesta combinação de Categoria Administrativa e Grande Região, oito cursos não receberam conceito ( SC ).
Os demais foram avaliados com conceito 1 ( oito cursos ), conceito 2 ( 127 cursos ), conceito 4 ( 69 cursos ) e conceito 5 ( 23 cursos ).
Entre os 33 cursos oferecidos por Instituições Públicas na região Sudeste, o conceito modal foi 4 ( 16 cursos ), e os demais cursos receberam os conceitos 1 ( um curso ), 2 ( seis cursos ), 3 ( quatro cursos ) e 5 ( seis cursos ).
Na região Sudeste, nesta categoria, nenhum curso ficou sem conceito ( SC ).
As Instituições Privadas concentraram 189 dos 224 cursos participantes da região Sul, 84,4 \% do total regional.
Desses, 99 ficaram com conceito 3, o conceito modal.
Nesta combinação de Categoria Administrativa e Grande Região, 34 cursos ficaram com conceito 4, 53 cursos obtiveram conceito 2, enquanto um único curso recebeu conceito 5.
Nenhum curso foi associado ao conceito 1, enquanto dois cursos ficaram sem conceito ( SC ).
As Instituições Públicas na região Sul participaram com 35 cursos ( 15,6 \% do total da região ), e o conceito modal foi 4, atribuído a 17 cursos.
Nove cursos foram avaliados com conceito 3 e cinco cursos, com conceito 5.
Um curso ficou sem conceito ( SC ), nenhum curso recebeu o conceito 1 e três cursos obtiveram conceito 2.
Na região Centro-Oeste, 92 dos 115 cursos participantes eram de Instituições Privadas ( 80,0 \% em termos regionais ).
Destes, 43 concentraram -se no conceito 2, conceito modal.
Dentre os demais, sete receberam conceito 1, um recebeu conceito 5 e seis receberam conceito 4.
Apenas um curso ficou sem conceito ( SC ) e 34 cursos obtiveram conceito 3.
Dos 23 cursos oferecidos por Instituições Públicas, sete foram avaliados com conceito 4, conceito modal, seguidos de seis cursos que obtiveram conceito 3.
Dos cursos restantes, um curso obteve o conceito 1, quatro cursos obtiveram conceito 5, e outros quatro, conceito 2.
Nesta região, apenas um curso de IES pública ficou sem conceito ( SC ).
O Gráfico 5.2 apresenta a distribuição do Conceito Enade segundo a Categoria Administrativa da IES.
Os cursos em IES Públicas ( linha azul ) apresentam uma poligonal mais à direita do que os em IES Privadas ( linha verde ) e, consequentemente, uma distribuição de Conceitos Enade com valores maiores. 5.3 CONCEITOS POR ORGANIZAÇÃO ACADÊMICA E POR GRANDE REGIÃO Na Tabela 5.3, encontra -se a distribuição dos conceitos atribuídos aos cursos participantes do Enade/2015 na Área de Ciências Contábeis, por Organização Acadêmica, segundo as Grandes Regiões brasileiras.
Dos 1.044 cursos de Ciências Contábeis participantes, 337 eram oferecidos em Universidades, 142 cursos na categoria Centros Universitários e 565 cursos oferecidos em Faculdades.
Esta distribuição corresponde, respectivamente, a 32,3 \%, 13,6 \% e 54,1 \% dos cursos.
De acordo com os dados apresentados, de todos os 43 cursos avaliados com conceito 5, 26 eram vinculados a Universidades.
Oito cursos oferecidos em Universidades ficaram sem conceito ( SC ).
Esse tipo de Organização Acadêmica teve o conceito 3 como modal, com 140 cursos ( 41,5 \% ).
Os demais cursos de Universidades avaliados receberam os conceitos 1 ( dois cursos ), 2 ( 66 cursos ), 4 ( 95 cursos ) e conceito 5 ( 26 cursos, como já mencionado ).
Entre os cursos em Centros Universitários, o conceito modal foi 3, com 61 cursos ( 43,0 \% ).
Neste tipo de Organização Acadêmica, apenas um curso ficou sem conceito ( SC ), e os outros cursos receberam os conceitos 1 ( três cursos ), 2 ( 45 cursos ), 4 ( 27 cursos ) e 5 ( cinco cursos ).
Nas Faculdades, sete dos 565 cursos ficaram sem conceito ( SC ), e 230 ( 40,7 \% ) ficaram alocados no conceito 3, conceito modal, seguido de 224 cursos que obtiveram conceito 2 e de 66 cursos que receberam conceito 4.
Os demais cursos neste tipo de Organização Acadêmica receberam os conceitos 1 ( 26 cursos ) e 5 ( 12 cursos ).
Considerando -se separadamente as regiões brasileiras, verifica -se que, na região Norte, as Universidades concentraram 13 dos 71 cursos participantes.
Nesta Organização Acadêmica, cinco cursos ficaram alocados ao conceito 4, conceito modal, seguidos de quatro cursos que obtiveram conceito 3, enquanto três cursos receberam conceito 2.
Nenhum curso ficou sem conceito ( SC ) ou tirou conceito 5 e apenas um curso obteve conceito 1.
Os Centros Universitários da região Norte foram representados por sete cursos.
Nenhum curso obteve conceito 1 ou 5 e nem ficou sem conceito ( SC ).
Apenas um curso obteve conceito 4 e os demais cursos receberam conceitos 2 e 3 ( três cursos cada ).
As Faculdades participaram com 51 cursos na região Norte, dos quais 23 receberam conceito 2 ( conceito modal ), quatro cursos obtiveram conceito 1, 17 cursos receberam conceito 3 e seis cursos obtiveram conceito 4.
Nenhum curso registrou conceito 5, e apenas um curso ficou sem conceito ( SC ).
Na região Nordeste, as Universidades participaram com 56 dos 190 cursos na Área de Ciências Contábeis.
O conceito modal foi 3, com 27 cursos, enquanto o conceito 2 foi atribuído a 15 cursos e o conceito 4, a 13 cursos.
Apenas um curso recebeu conceito 5 e nenhum curso obteve conceito 1 ou ficou sem conceito ( SC ).
Os Centros Universitários contaram com 11 cursos participantes na região Nordeste.
Desta Organização Acadêmica, nenhum curso ficou sem conceito ( SC ) ou tirou conceito 1, e oito cursos obtiveram conceito 3, conceito modal.
Aos três cursos restantes foram atribuídos os conceitos 2, 4 ou 5.
As Faculdades foram representadas por 123 cursos na região Nordeste, dos quais 54 receberam conceito 2, conceito modal.
Dos demais cursos, 44 tiveram conceito 3, seguidos de 13 cursos com conceito 4, nove cursos com conceito 1 e um curso com conceito 5.
Nessa região, apenas dois cursos ficaram sem conceito ( SC ).
Na região Sudeste, as Universidades concentraram 139 dos 444 cursos da região.
Entre os cursos em Universidades, nessa região, o conceito modal foi 3, com 53 cursos, enquanto nenhum curso ficou com conceito 1 e seis cursos ficaram sem conceito ( SC ).
Os demais cursos receberam os conceitos 2 ( 27 cursos ), 4 ( 38 cursos ), 5 ( 15 cursos ).
Os Centros Universitários participaram com 85 cursos na região Sudeste, dos quais 33 obtiveram o conceito 3 ( conceito modal ), 30 cursos, conceito 2, e 15 cursos, conceito 4.
Três cursos obtiveram conceito 1, e outros três, conceito 5.
Apenas um curso nessa região ficou sem conceito ( SC ).
As Faculdades foram representadas por 220 cursos na região Sudeste, dos quais seis receberam o conceito 1 e 11 cursos receberam conceito 5.
Os demais se \textbf{distribuíram} nos conceitos 2 ( 76 cursos ), 3 ( 94 cursos, conceito modal ) e 4 ( 32 cursos ).
Apenas um curso ficou sem conceito ( SC ) nessa categoria e região.
Dos 224 cursos da região Sul, 98 eram de Universidades, para os quais o conceito modal foi 3, atribuído a 42 cursos.
Nesse tipo de organização, 2 cursos ficaram sem conceito ( SC ), e nenhum curso obteve conceito 1.
Os demais cursos receberam os conceitos 5 ( 6 cursos ), 2 ( 16 cursos ), 4 ( 32 cursos ).
Dos 25 cursos participantes de Centros Universitários da região Sul, nenhum ficou sem conceito ( SC ) ou obteve conceitos 1 ou 5.
Cinco cursos receberam conceito 2, e os conceitos 3 e 4 foram obtidos por dez cursos cada.
Com relação às Faculdades, 101 cursos foram avaliados, sendo que um curso ficou sem conceito ( SC ) e nenhum curso obteve conceitos 1 ou 5.
Os demais cursos receberam conceitos 4 ( nove cursos ), 2 ( 35 cursos ) e 3, o conceito modal, ( 56 cursos ).
Na região Centro-Oeste, 31 dos 115 cursos eram de Universidades.
Nesse tipo de organização, nenhum curso ficou sem conceito ( SC ), e apenas um curso obteve conceito 1.
O conceito modal foi 3, atribuído a 14 cursos, seguido do conceito 4, obtido por sete cursos.
Cinco cursos receberam conceito 2 e quatro cursos foram alocados no conceito 5.
Dos 14 Centros Universitários da região Centro-Oeste, nenhum ficou sem conceito ( SC ) ou obteve conceitos 1 ou 4.
Apenas um curso obteve conceito 5, e sete receberam conceito 3, o conceito modal, seguidos de 6 cursos que obtiveram o conceito 2.
As Faculdades foram representadas por 70 cursos, sendo que dois cursos ficaram sem conceito ( SC ) e nenhum obteve conceito 5.
Quanto aos demais cursos, seis foram alocados no conceito 4, sete no conceito 1, 19 cursos no conceito 3 e 36 no conceito 2 ( conceito modal ).
O Gráfico 5.3 apresenta a distribuição do Conceito Enade segundo a Organização Acadêmica da IES.
Os cursos em Universidades ( linha azul ) apresentam uma poligonal mais à direita do que as demais, denotando uma melhor distribuição de Conceitos Enade.
No outro extremo, temos os cursos oferecidos em Faculdades ( linha vermelha ) com a poligonal mais à esquerda e, consequentemente, uma distribuição de Conceitos Enade com menores valores.
IES de Centros Universitários ( linha verde ) apresentam uma situação intermediária.

% matched lemmas: distribuir
\end{textsample}
